\documentclass[11pt,a4paper]{article}
\usepackage[margin=25mm]{geometry}
\usepackage{amsmath,amssymb,amsthm,hyperref}
\setlength{\emergencystretch}{2em}
\newtheorem{lemma}{Lemma}
\newtheorem{theorem}[lemma]{Theorem}
\newcommand{\gen}[1]{\langle#1\rangle}
\title{Transitive label budgets on six residual fixed points}
\author{Human Compression project: L044 continuation}
\date{}
\begin{document}\maketitle
Use the critical singleton-end setting of the accepted residual-deficit
proof: $G=\gen{\rho_0,\ldots,\rho_{r-1}}\leq A_n$ is transitive,
is a connected string C-group, and has every maximal omission intransitive;
$G_0$ has orbits $\{a\}$ and its complement;
$\rho_0=(a\ b)\delta$, $K=G_{>1}$, $R=G_{>2}$ and $\sigma=\rho_1$.
The odd involution $\delta$ has $t$ pairs, with union $P$, and an even
fixed set $F$ of size $f$ in $\Omega\setminus\{a,b\}$.
Assume $n=2r\geq14$, so $K$ has $d=t+f/2-1$ labels.
All block-image omission-orbit statements below use the complete
accepted actual-deficit proof. Whiston's bound and the complete
odd-degree maximal-string theorem, including degree five, remain paid.

\begin{lemma}[Single-transposition restrictions]\label{single}
If $K$ is transitive on $F$, $f\geq4$, and all nonidentity distinguished
restrictions to $F$ are single transpositions, they are exactly the
$f-1$ successive edges of a natural path, each with one label.
\end{lemma}
\begin{proof}
Their edge graph is connected. Three distinct edges incident to one
vertex would have pairwise noncommuting labels, impossible in a string.
A cycle is likewise impossible: labels chosen on its distinct edges
would require its line graph, a cycle, to embed in the path of label
positions. Hence the simple edge graph is a path.

Choose a label on each edge of that path. Their positions form a path
in the integer line, so are consecutive in one of the two orders.
Any repeated label on an interior edge would have to be adjacent to
both chosen neighboring labels, whose unique common neighbor is the
already occupied position. For an end edge, its neighboring edge is
interior because $f\geq4$; both positions adjacent to that neighboring
label are already occupied. A second end-edge label is also impossible.
\end{proof}

\begin{lemma}[Minimum transitive fixed-set budget]\label{budget}
Let $E$ be a set of $h$ distinguished $K$ labels whose restrictions
are transitive on $F$. Then $h\geq f/2$.
If equality holds and $t\geq3$ is odd, the other $t-1$ pair images
generate the ordinary $S_t$ path, at consecutive positions $J$,
and every label of $E$ has trivial pair image.
\end{lemma}
\begin{proof}
Deleting a label outside $E$ leaves the fixed-set action transitive.
Its pair image cannot be generated by the other pair images, as that
would leave the full block-orbit partition unchanged. In particular
the $d-h$ pair images outside $E$ are independent. Whiston gives
$d-h\leq t-1$, or $h\geq f/2$.

At equality those $t-1$ images generate $S_t$ and form its ordinary
path by the accepted odd-degree maximal-string theorem. For $t=3$
this just says that two independent involutions in $S_3$ are distinct
transpositions. Consecutive noncommutations force consecutive positions
in the original string. Omitting any path mark leaves two pair orbits.
If an $E$ label joined them, the full omission would be transitive both
on pairs and on $F$, contradicting the actual orbit deficit. Thus every
$E$ label preserves every path cut; their common stabilizer fixes every
pair-block. This proves the last assertion.
\end{proof}

\begin{theorem}[The three-label six-point branch]\label{three}
Suppose $f=6$ and three distinguished $K$ restrictions are transitive
on $F$. Then necessarily $t=3$, $n=14$, $r=7$, and $\rho_2$ is one
of those three labels. In particular the case $t\geq5$ is impossible.
\end{theorem}
\begin{proof}
Let $E$ be the three labels. By Lemma~\ref{budget} their pair images
are trivial, while $J$ is an ordinary path on the $t$ pairs.
Here $t\geq3$ because $n=2t+8\geq14$.

First suppose $\rho_2\in J$. Then $J$ starts at index2 and all three
extras follow it. The group $R$ retains $E$, so $F$ is an $R$-orbit
of size six. On pairs, $R$ fixes the first endpoint pair and is
transitive on the other $t-1$. Its point orbits in $P$ therefore have
sizes at most two, $t-1$, or $2(t-1)$: in a transitive pair-block
orbit, a point orbit meets every block in the same number of points.
The only possible size-six collision for odd $t$ is $t=7$, with
a row meeting each of the six pair-blocks once. Every $E$ label has
trivial pair image and preserves that $R$-orbit, so fixes it pointwise.
It cannot be $\sigma(F)$, since equivariance would then make all of
$E$ fix $F$, whereas $E$ is transitive on $F$. In every case $\sigma$
must preserve $F$. So does $K$, contradicting $G_0$ transitivity.

Thus $\rho_2\in E$. Now $R$ contains all of $J$, so its point
orbits in $P$ have size $t$ or $2t$. If $t\geq7$, their sizes
exceed six, and $P$ is their invariant union under $\sigma$.
This again contradicts $G_0$ transitivity. If $t=5$, the only
possible collision is a five-point row $O\subseteq P$ and an
$R$-orbit $U\subseteq F$ of size five. The two surviving $E$
labels have trivial pair images, and hence fix $O$ pointwise.
If $\sigma$ exchanged $O,U$, they would fix $U$ pointwise as well,
and therefore also fix its single-point complement in $F$.
Then the entire $E$ action on $F$ would be generated by the one
omitted involution, which is not transitive on six points.
So that collision cannot occur; $\sigma$ preserves $P$, again
impossible. Only $t=3$ remains.
\end{proof}

\paragraph{Literal six-point restrictions.}
For a transitive action on six fixed points, a distinguished triple
transposition supplies three disjoint edges. Extend them to a spanning
tree of the full fixed-set edge graph; just two further edges are needed,
so at most three distinguished labels are transitive on $F$.
Lemma~\ref{budget} makes the minimum exactly three, and
Theorem~\ref{three} applies. A double transposition similarly gives a
transitive set of at most four labels. If all restrictions are single
transpositions, Lemma~\ref{single} puts the action within the separately
proved natural six-point-path exclusion.

\paragraph{Attribution and boundary.}
The source is Cameron--Fernandes--Leemans--Mixer,
\emph{Highest rank of a polytope for $A_n$}, arXiv:1605.09173v1,
physical pages38--44. The path-label observation extends the separate
accepted four-fixed proof by \texttt{rewrite\_termination}; the
transitive-label budget and the present six-point reductions use that
proof's actual-omission and equivariant-orbit approach. The exact
natural six-point exclusion is a separate new bounded dependency,
not part of Lemma~\ref{single} alone.

The $n=14$ three-label boundary is not excluded here. Nor is the
four-label six-point case, an intransitive action on six fixed points,
or any general larger fixed set. The theorem performs no finite
enumeration and invokes no projected IP, alternating induction,
whole-source recovery, new original credit or conjecture credit.
Distinct review remains pending.
\end{document}
